\ApiScopeSource{api-path-geo-planar}{semio / viz / geo / planar}
\SemioTableLong[text-size=8pt]{\ApiTitle{api-path-geo-planar}}{0.20,0.13,0.15,0.52}{\ApiKey & \ApiType & \ApiDefault & \ApiMeaning}{%
  \SemioTableRow{\Key{width} & \Key{number} & \Key{76} & \ApiText{Maximum fitted display width in millimetres. Raw planar coordinates are scaled uniformly.}{Maximale angepasste Darstellungsbreite in Millimetern. Rohe planare Koordinaten werden gleichmäßig skaliert.}}
  \SemioTableRow{\Key{height} & \Key{number} & \Key{44} & \ApiText{Maximum fitted display height in millimetres. Raw planar coordinates are scaled uniformly.}{Maximale angepasste Darstellungshöhe in Millimetern. Rohe planare Koordinaten werden gleichmäßig skaliert.}}
  \SemioTableRow{\Key{palette} & \Key{string} & \Key{sequential} & \ApiText{Colour ramp used for binned values or density levels.}{Farbverlauf für gebündelte Werte oder Dichteniveaus.}}
  \SemioTableRow{\Key{stroke} & \Key{string} & \Key{semio-chrome-border-emphasized} & \ApiText{Hexagon outline colour as a native colour expression. Density and heat outlines use their palette colour.}{Umrissfarbe der Sechsecke als nativer Farbausdruck. Dichte- und Wärmeumrisse nutzen ihre Palettenfarbe.}}
  \SemioTableRow{\Key{lineWidth} & \Key{string} & \Key{semio@stroke@hairline} & \ApiText{Hexagon and density outline width as a TeX dimension from the theme hairline token.}{Umrissbreite der Sechsecke und Dichteflächen als TeX-Dimension aus dem Haarlinienwert des Themas.}}
  \SemioTableRow{\Key{opacity} & \Key{number} & \Key{1} & \ApiText{Drawing opacity; hex fill uses the same value, heat fill uses 0.3 times it.}{Zeichendeckkraft; Sechseckfüllungen nutzen denselben Wert, Wärmefüllungen das 0,3-Fache.}}
  \SemioTableRow{\Key{paletteSteps} & \Key{integer} & \Key{5} & \ApiText{Positive number of final colour-ramp levels. One level uses the palette midpoint.}{Positive Anzahl endgültiger Farbverlaufsstufen. Eine Stufe nutzt die Palettenmitte.}}
}
